% GENERATED FILE. Edit canonical lessons, then run npm run generate:books.
% canonical-source-hash: fnv1a64:97afcbc4e9fb44cd
% canonical-lessons: ML-C567-anubhavikkuka, ML-C567-nerituka, ML-C567-piticcunilkkuka, ML-C567-marikatakkuka, ML-C567-poruttappetuka

\chapter{Facing Hard Times}
\label{ch:ml-facing-hard-times}
\hlchaptermodality{567}

\begin{chapteropening}
\textbf{By the end of this chapter:} \emph{I can say to go through, to confront, to hold out, to overcome and to adapt.}

Complete the last lesson of chapter 567: I can say to go through, to confront, to hold out, to overcome and to adapt.
\end{chapteropening}

\section[\texorpdfstring{\ml{അനുഭവിക്കുക} (anubhavikkuka)}{അനുഭവിക്കുക (anubhavikkuka)}]{\ml{അനുഭവിക്കുക} (anubhavikkuka) — to go through; to experience}

\label{lesson:ML-C567-anubhavikkuka}



\begin{quote}
Before the new one: say the Malayalam for gardening, then the Malayalam for to meditate.
\end{quote}

\subsection*{You'll want to know: \ml{അനുഭവിക്കുക}}
\textbf{\ml{അനുഭവിക്കുക}} — \emph{anubhavikkuka} — \textquotedblleft{}to go through; to experience\textquotedblright{}.

\begin{quote}
\ml{കുട്ടിക്കാലത്ത്} \ml{അവർ} \ml{ഒരുപാട്} \ml{കഷ്ടപ്പാട്} \ml{അനുഭവിച്ചു}.
\end{quote}

\begin{quote}
\emph{kuṭṭikkālattŭ avar orupāṭŭ kaṣṭappāṭŭ anubhaviccu.}
\end{quote}

\begin{quote}
They went through a lot of hardship in their childhood.
\end{quote}

Past: \ml{അനുഭവിച്ചു}, \emph{anubhaviccu}.

The noun is \ml{അനുഭവം}, a lived experience: \ml{എനിക്ക്} \ml{ഒരു} \ml{നല്ല} \ml{അനുഭവം} \ml{ഉണ്ടായി}.

\begin{cousinweb}
\textbf{\ml{അനുഭവിക്കുക}} (\emph{anubhavikkuka}) is built from Sanskrit \emph{anubhava}, \textquotedblleft{}perception, fruition, understanding\textquotedblright{} (Monier-Williams), with the Malayalam verb ending -\ml{ിക്കുക}.
\end{cousinweb}

\begin{grammarlens}[title={What you've built}]
One more word for describing what you lived and felt.
\end{grammarlens}

\subsection*{Your turn}
\begin{itemize}
\raggedright
  \item \emph{Say it:} \emph{anubhavikkuka}
  \item \emph{Say it:} \emph{anubhavikkuka}, once more
  \item \emph{Say it:} \emph{kuṭṭikkālattŭ avar orupāṭŭ kaṣṭappāṭŭ anubhaviccu.} — They went through a lot of hardship in their childhood.
  \item \emph{Recall:} say the Malayalam for gardening, then the Malayalam for to meditate, then say \emph{anubhavikkuka} again
\end{itemize}

\begin{tcolorbox}[breakable,colback=teal!4,colframe=teal!35!black,title={Before you move on}]
What does \ml{അനുഭവിക്കുക} mean? (\textquotedblleft{}To go through; to experience\textquotedblright{}.) Say it once more.
\end{tcolorbox}
\section[\texorpdfstring{\ml{നേരിടുക} (nēriṭuka)}{നേരിടുക (nēriṭuka)}]{\ml{നേരിടുക} (nēriṭuka) — to confront; to face}

\label{lesson:ML-C567-nerituka}



\begin{quote}
Before the new one: say the Malayalam for to meditate, then the Malayalam for to go through.
\end{quote}

\subsection*{You'll want to know: \ml{നേരിടുക}}
\textbf{\ml{നേരിടുക}} — \emph{nēriṭuka} — \textquotedblleft{}to confront; to face\textquotedblright{}.

\begin{quote}
\ml{ആ} \ml{കുടുംബം} \ml{പ്രശ്നങ്ങളെ} \ml{ധൈര്യത്തോടെ} \ml{നേരിട്ടു}.
\end{quote}

\begin{quote}
\emph{ā kuṭumbaṁ praśnaṅṅaḷe dhairyattōṭe nēriṭṭu.}
\end{quote}

\begin{quote}
That family faced its problems with courage.
\end{quote}

Past: \ml{നേരിട്ടു}, \emph{nēriṭṭu}.

\ml{നേർ} facing + \ml{ഇടുക}; \ml{നേരിട്ട്}, directly, is the same form used as an adverb.

\begin{cousinweb}
\textbf{\ml{നേരിടുക}} (\emph{nēriṭuka}) is native Dravidian: DEDR 3770 lists \ml{നേരിടുക} itself, \textquotedblleft{}to face, to meet, to oppose\textquotedblright{}, built on \ml{നേർ}, \textquotedblleft{}what is even, facing\textquotedblright{} (DEDR 3770).
\end{cousinweb}

\begin{grammarlens}[title={What you've built}]
One more word for describing what you lived and felt.
\end{grammarlens}

\subsection*{Your turn}
\begin{itemize}
\raggedright
  \item \emph{Say it:} \emph{nēriṭuka}
  \item \emph{Say it:} \emph{nēriṭuka}, once more
  \item \emph{Say it:} \emph{ā kuṭumbaṁ praśnaṅṅaḷe dhairyattōṭe nēriṭṭu.} — That family faced its problems with courage.
  \item \emph{Recall:} say the Malayalam for to meditate, then the Malayalam for to go through, then say \emph{nēriṭuka} again
\end{itemize}

\begin{tcolorbox}[breakable,colback=teal!4,colframe=teal!35!black,title={Before you move on}]
What does \ml{നേരിടുക} mean? (\textquotedblleft{}To confront; to face\textquotedblright{}.) Say it once more.
\end{tcolorbox}
\section[\texorpdfstring{\ml{പിടിച്ചുനിൽക്കുക} (piṭiccunilkkuka)}{പിടിച്ചുനിൽക്കുക (piṭiccunilkkuka)}]{\ml{പിടിച്ചുനിൽക്കുക} (piṭiccunilkkuka) — to hold out; to cope}

\label{lesson:ML-C567-piticcunilkkuka}



\begin{quote}
Before the new one: say the Malayalam for to go through, then the Malayalam for to confront.
\end{quote}

\subsection*{You'll want to know: \ml{പിടിച്ചുനിൽക്കുക}}
\textbf{\ml{പിടിച്ചുനിൽക്കുക}} — \emph{piṭiccunilkkuka} — \textquotedblleft{}to hold out; to cope\textquotedblright{}.

\begin{quote}
\ml{ജോലി} \ml{പോയിട്ടും} \ml{ആ} \ml{കുടുംബം} \ml{പിടിച്ചുനിന്നു}.
\end{quote}

\begin{quote}
\emph{jōli pōyiṭṭuṁ ā kuṭumbaṁ piṭiccuninnu.}
\end{quote}

\begin{quote}
Even after losing the job, that family held out.
\end{quote}

Past: \ml{പിടിച്ചുനിന്നു}, \emph{piṭiccuninnu}.

\ml{പിടിച്ച്}, holding on, + \ml{നിൽക്കുക}, to stand: to stand firm.

\begin{cousinweb}
\textbf{\ml{പിടിച്ചുനിൽക്കുക}} (\emph{piṭiccunilkkuka}) is native Dravidian, a compound of \ml{പിടിക്കുക}, \textquotedblleft{}to hold\textquotedblright{} (DEDR 4148), and \ml{നിൽക്കുക}, \textquotedblleft{}to stand\textquotedblright{} (DEDR 3675).
\end{cousinweb}

\begin{grammarlens}[title={What you've built}]
One more word for describing what you lived and felt.
\end{grammarlens}

\subsection*{Your turn}
\begin{itemize}
\raggedright
  \item \emph{Say it:} \emph{piṭiccunilkkuka}
  \item \emph{Say it:} \emph{piṭiccunilkkuka}, once more
  \item \emph{Say it:} \emph{jōli pōyiṭṭuṁ ā kuṭumbaṁ piṭiccuninnu.} — Even after losing the job, that family held out.
  \item \emph{Recall:} say the Malayalam for to go through, then the Malayalam for to confront, then say \emph{piṭiccunilkkuka} again
\end{itemize}

\begin{tcolorbox}[breakable,colback=teal!4,colframe=teal!35!black,title={Before you move on}]
What does \ml{പിടിച്ചുനിൽക്കുക} mean? (\textquotedblleft{}To hold out; to cope\textquotedblright{}.) Say it once more.
\end{tcolorbox}
\section[\texorpdfstring{\ml{മറികടക്കുക} (maṟikaṭakkuka)}{മറികടക്കുക (maṟikaṭakkuka)}]{\ml{മറികടക്കുക} (maṟikaṭakkuka) — to overcome; to get past}

\label{lesson:ML-C567-marikatakkuka}



\begin{quote}
Before the new one: say the Malayalam for to confront, then the Malayalam for to hold out.
\end{quote}

\subsection*{You'll want to know: \ml{മറികടക്കുക}}
\textbf{\ml{മറികടക്കുക}} — \emph{maṟikaṭakkuka} — \textquotedblleft{}to overcome; to get past\textquotedblright{}.

\begin{quote}
\ml{അസുഖം} \ml{മറികടന്ന്} \ml{അവൻ} \ml{പരീക്ഷ} \ml{എഴുതി}.
\end{quote}

\begin{quote}
\emph{asukhaṁ maṟikaṭannŭ avan parīkṣa eḻuti.}
\end{quote}

\begin{quote}
He got past his illness and sat the exam.
\end{quote}

Past: \ml{മറികടന്നു}, \emph{maṟikaṭannu}.

\ml{മറി} a turn, a turning over + \ml{കടക്കുക} to cross, already familiar.

\begin{cousinweb}
\textbf{\ml{മറികടക്കുക}} (\emph{maṟikaṭakkuka}) is native Dravidian, a compound of \ml{മറി}, \textquotedblleft{}a turn, a turning over\textquotedblright{} (DEDR 4761), and \ml{കടക്കുക}, \textquotedblleft{}to pass over, to surpass\textquotedblright{} (DEDR 1109).
\end{cousinweb}

\begin{grammarlens}[title={What you've built}]
One more word for describing what you lived and felt.
\end{grammarlens}

\subsection*{Your turn}
\begin{itemize}
\raggedright
  \item \emph{Say it:} \emph{maṟikaṭakkuka}
  \item \emph{Say it:} \emph{maṟikaṭakkuka}, once more
  \item \emph{Say it:} \emph{asukhaṁ maṟikaṭannŭ avan parīkṣa eḻuti.} — He got past his illness and sat the exam.
  \item \emph{Recall:} say the Malayalam for to confront, then the Malayalam for to hold out, then say \emph{maṟikaṭakkuka} again
\end{itemize}

\begin{tcolorbox}[breakable,colback=teal!4,colframe=teal!35!black,title={Before you move on}]
What does \ml{മറികടക്കുക} mean? (\textquotedblleft{}To overcome; to get past\textquotedblright{}.) Say it once more.
\end{tcolorbox}
\section[\texorpdfstring{\ml{പൊരുത്തപ്പെടുക} (poruttappeṭuka)}{പൊരുത്തപ്പെടുക (poruttappeṭuka)}]{\ml{പൊരുത്തപ്പെടുക} (poruttappeṭuka) — to adapt; to adjust}

\label{lesson:ML-C567-poruttappetuka}



\begin{quote}
Before the new one: say the Malayalam for to hold out, then the Malayalam for to overcome.
\end{quote}

\subsection*{You'll want to know: \ml{പൊരുത്തപ്പെടുക}}
\textbf{\ml{പൊരുത്തപ്പെടുക}} — \emph{poruttappeṭuka} — \textquotedblleft{}to adapt; to adjust\textquotedblright{}.

\begin{quote}
\ml{പുതിയ} \ml{നഗരവുമായി} \ml{പൊരുത്തപ്പെടാൻ} \ml{കുറച്ച്} \ml{സമയം} \ml{എടുത്തു}.
\end{quote}

\begin{quote}
\emph{putiya nagaravumāyi poruttappeṭān kuṟaccŭ samayaṁ eṭuttu.}
\end{quote}

\begin{quote}
It took a little time to adjust to the new city.
\end{quote}

Past: \ml{പൊരുത്തപ്പെട്ടു}, \emph{poruttappeṭṭu}.

\ml{പൊരുത്തം} a fit, a match + \ml{പെടുക}. What one adapts to takes -\ml{ഉമായി}.

\begin{cousinweb}
\textbf{\ml{പൊരുത്തപ്പെടുക}} (\emph{poruttappeṭuka}) is native Dravidian, built from \ml{പൊരുത്തം}, \textquotedblleft{}suitableness, accord\textquotedblright{} (DEDR 4541), and the helper verb \ml{പെടുക}.
\end{cousinweb}

\begin{grammarlens}[title={What you've built}]
One more word for describing what you lived and felt.
\end{grammarlens}

\subsection*{Your turn}
\begin{itemize}
\raggedright
  \item \emph{Say it:} \emph{poruttappeṭuka}
  \item \emph{Say it:} \emph{poruttappeṭuka}, once more
  \item \emph{Say it:} \emph{putiya nagaravumāyi poruttappeṭān kuṟaccŭ samayaṁ eṭuttu.} — It took a little time to adjust to the new city.
  \item \emph{Recall:} say the Malayalam for to hold out, then the Malayalam for to overcome, then say \emph{poruttappeṭuka} again
\end{itemize}

\begin{tcolorbox}[breakable,colback=teal!4,colframe=teal!35!black,title={Before you move on}]
What does \ml{പൊരുത്തപ്പെടുക} mean? (\textquotedblleft{}To adapt; to adjust\textquotedblright{}.) Say it once more.
\end{tcolorbox}
